%% =====================================================================
%%  Quantum Electronic Nose — Master Document ("mini-thesis")
%%
%%  A single living document that collects everything that has been done,
%%  is being done, and is planned on the quantum e-nose project. Sections
%%  of this file are the source material for:
%%
%%    - the SISPAD 2026 abstract (papers/IEEE-conference-template-062824/
%%      and papers/sispad2026_abstract.tex),
%%    - the full IEEE TED journal paper,
%%    - internal documentation and onboarding notes,
%%    - the roadmap for taking the broader project forward.
%%
%%  Compile with pdflatex. Figures are placeholder frames; drop in real
%%  PDFs under fig/ when they are ready.
%% =====================================================================
\documentclass[11pt,a4paper,onecolumn]{article}

\usepackage[margin=1in]{geometry}
\usepackage{amsmath,amssymb,amsfonts}
\usepackage{graphicx}
\usepackage{xcolor}
\usepackage{hyperref}
\usepackage{booktabs}
\usepackage{enumitem}
\usepackage{microtype}
\usepackage{titlesec}
\usepackage{siunitx}
\hypersetup{colorlinks=true,linkcolor=blue!40!black,citecolor=blue!40!black,urlcolor=blue!60!black}

\titleformat*{\section}{\Large\bfseries}
\titleformat*{\subsection}{\large\bfseries}

\newcommand{\code}[1]{\texttt{\small #1}}
\newcommand{\TODO}[1]{{\color{red}\textbf{TODO:} #1}}
\newcommand{\placeholder}[2]{%
  \fbox{\parbox[c][#2][c]{\dimexpr#1-2\fboxsep-2\fboxrule\relax}{\centering\textit{placeholder figure}}}%
}

\title{\Huge A Room-Temperature Quantum Electronic Nose\\[4pt]
\Large Rank-1 Projected NEGF--SCBA Simulation of\\
BEOL-Compatible Oxide Resonant Tunneling Diodes\\
for Inelastic-Tunneling Molecular Discrimination\\[8pt]
\large\textit{Master project document — internal, living}}
\author{Abhineet Agarwal\\
\small with the Quantum E-Nose group\\
\small assisted by Claude (Anthropic)}
\date{Last updated: \today}

\begin{document}
\maketitle

\begin{abstract}\noindent
This document is the authoritative, living record of the quantum
electronic nose project. It states the scientific problem, reviews prior
art in the group's lineage (Lambe--Jaklevic, Patil, Pandey), develops the
complete formalism we use---including the novel rank-1 projected
self-consistent Born approximation (SCBA), converged via Anderson mixing
(DIIS), and the closed-form analytic $d^2I/dV^2$ (available but not used
in production; the self-consistent SCBA current is smooth enough for raw
numerical double-differentiation)---documents the material-stack selection
that led to the BEOL-compatible
ZnO/Mg\textsubscript{0.3}Zn\textsubscript{0.7}O primary stack, collects
validation evidence against the Patil 1D benchmark (0.55\% at NDR peak),
presents SCBA production results on synthetic single- and two-mode
analytes (Baseline, Mol\_A, Mol\_B, Mol\_AB) at 10--300~K, and sets out
a prioritised roadmap of open items. It is written so that distinct
subsections can be lifted with minimal editing into the SISPAD conference
abstract, into an IEEE~TED journal manuscript, and into internal
onboarding documentation. Every claim is cross-referenced to the code
that produces it.
\end{abstract}

\tableofcontents
\newpage

% =====================================================================
\section{Introduction and Motivation}
% =====================================================================

\subsection{The problem: 300\,K vibrational fingerprinting is hard}
Molecular olfaction ultimately discriminates species by their vibrational
spectra. Inelastic Electron Tunneling Spectroscopy (IETS), discovered by
Jaklevic and Lambe in 1966 in metal--insulator--metal (MIM) junctions, is
the direct electrical analogue: electrons tunnelling across a thin
insulator excite molecular vibrations in the barrier, opening inelastic
channels that leave narrow features in $d^2I/dV^2$ at
$eV = \hbar\omega_\nu$. The method has unmatched chemical specificity but
historically requires cryogenic operation ($T \lesssim 10$~K) because the
Fermi-level thermal broadening $\sim\! 3k_BT$ at 300~K is already
$\sim\! 78$~meV, wider than the spacing between many IETS peaks. This
thermal ceiling is the single biggest obstacle to using IETS as a room
temperature chemical sensor.

\subsection{The PI's idea: a quantum filter beats thermal broadening}
The core idea of this project is to put a quantum energy filter
\emph{in series with the molecular junction}, so that the energy window
seen by the vibrational channel is set by the filter's spectral width
rather than by the contact Fermi function. A double-barrier resonant
tunneling diode (RTD) is a natural filter: its resonant level has an
intrinsic width
$\Gamma_{\mathrm{res}} \!\ll\! k_BT$ and acts as a \emph{sharper than
thermal} band-pass that preserves vibrational peaks at 300~K. The PI's
hypothesis is that the combination RTD + embedded molecule turns IETS into
a room-temperature sensing modality. Quantum dots are an alternative
filter; RTDs are chosen here because they are planar-processable, area-
scalable, and already demonstrated in multiple oxide material systems
compatible with back-end-of-line (BEOL) CMOS integration.

\subsection{Lineage: Patil, Pandey, and the group}
\begin{itemize}[leftmargin=1.3em]
\item \textbf{Patil \textit{et al.}}, \textit{Sci.\ Rep.}~2018 (``A quantum
biomimetic electronic nose sensor'') is the direct antecedent. Patil
simulated a 1D GaAs/AlGaAs RTD with SCBA and demonstrated that a narrow
RTD resonance does preserve IETS fingerprints at 300~K in a toy 1D model
with a point-like Einstein vibron coupled by $D_0^2 = 0.1$~eV$^2$. The
matlab source (\code{rtd2modes\_1d.m}) is preserved in this repo as the
reference implementation.
\item \textbf{Nidhi Pandey} (prior PhD in the group) extended the RTD-IETS
picture along two independent axes: (a) a demonstration that
enantioselectivity is in principle possible within this framework, and
(b) a graph-based analysis linking a molecule's vibrational spectrum to
its olfactory percept. These two directions are relevant future
integration points for this project.
\item \textbf{Broader group activity} includes instrumentation work on
cryogenic IETS and a prior cryogenic IETS-tongue demonstration. Those
provide experimental grounding but do not address the 300~K problem.
\end{itemize}

\subsection{This project's scope: 3D, realistic stack, molecular
discrimination}
The contribution of this project is to take the Patil idea out of its 1D
toy setting and answer four questions quantitatively:
\begin{enumerate}[leftmargin=1.3em,itemsep=2pt]
\item \textbf{Quasi-3D.} How does the quasi-3D structure of a real RTD
device change the IETS response relative to Patil's 1D model? The key
technical challenge is handling the transverse-mode structure without
exploding the computational cost of the SCBA.
\item \textbf{BEOL stack.} Which material system in the low conduction-band
offset ($\Delta E_c \!\approx\! 0.5$~eV), BEOL-compatible oxide family can
support a working RTD with NDR and operate in the molecular vibrational
energy window?
\item \textbf{Discrimination.} Given two synthetic analytes with known
vibrational spectra, can the device discriminate them at 300~K at a
reviewer-defensible signal-to-noise ratio?
\item \textbf{Mixtures.} Does the IETS fingerprint of a mixed analyte
behave linearly, so that standard multivariate pattern recognition applies?
\end{enumerate}

\subsection{What is explicitly \emph{not} the focus right now}
To keep the project tractable and the claims technically sound, a number
of adjacent problems are deliberately deferred:
\begin{itemize}[leftmargin=1.3em]
\item \textbf{Realistic molecules.} The analytes used for discrimination
are synthetic (\code{Mol\_A/B/AB}) dummies with one or two Einstein modes.
Real DFT-derived VOC spectra (benzene, thiophene, naphthalene, dimethyl
sulfide, acetone, etc.) are already in the molecular database but are
deferred until the methodology paper is submitted; they add a confound
between molecule-geometry and IETS-signature.
\item \textbf{Self-consistent Poisson.} The bias profile is currently a
linear drop. A self-consistent Poisson loop on $U_{\mathrm{bias}}$ is a
planned follow-up.
\item \textbf{Arbitrary-stack / external-simulator interop.} Out of scope.
\item \textbf{Full image-charge renormalisation} of the vibron frequency
by the oxide dielectric environment ($\sim 50$~meV shift). Noted but not
corrected since the analytes are synthetic.
\end{itemize}

% =====================================================================
\section{Prior Art and Lineage}
% =====================================================================

\subsection{Lambe--Jaklevic MIM IETS (1966)}
Original demonstration of IETS in Al/Al$_2$O$_3$/Pb junctions. Establishes
the $eV = \hbar\omega$ peak-in-$d^2I/dV^2$ signature, the thermal broadening
ceiling, and the central role of barrier quality. The cryogenic ceiling is
the gap this project is trying to close.

\subsection{Patil \textit{et al.}, Sci.\ Rep.\ 2018}
Core paper. Uses a 1D SCBA NEGF to show that a double-barrier RTD
resonance acts as a narrow energy window that preserves IETS peaks at 300~K.
Key numerical parameters used as ground truth here:
symmetric bias convention $\mu_{L,R}(V) = E_F \pm V/2$; $E_F = 20$~meV;
electron--phonon coupling $D_0^2 = 0.1$~eV$^2$; bulk phonon off (their
Dnu\,=\,0 default); GaAs/AlAs-like barriers. The reference implementation is
\code{papers/rtd2modes\_1d.m}. We use it as our Tier-1 $\gamma$ benchmark
(see \S\ref{sec:validation}).

\subsection{Pandey (group PhD) --- enantioselectivity and spectrum--percept
mapping}
Two results from Pandey's thesis work are on the follow-up roadmap:
(i)~the RTD-IETS transport integral is sensitive to the chirality-dependent
projection of the molecular vibration onto the transport direction, so in
principle enantiomers produce distinct $d^2I/dV^2$ signatures;
(ii)~a graph-neural-net analysis mapping vibrational spectra to olfactory
percepts. Both are logical extensions of the present simulator once the
methodology paper is out.

\subsection{Galperin--Ratner--Nitzan; Troisi--Ratner}
Reviews the inelastic transport formalism (Galperin \textit{et al.}, JPCM
2007) and propensity rules for IETS selection (Troisi \& Ratner, \textit{Nano
Lett.}\ 2006). These set the standard against which our SCBA
implementation is checked.

\subsection{Krishnamoorthy \textit{et al.}, Solid-State Electron.\ 2002 --- ZnO/MgZnO RTD}
First experimentally demonstrated resonant tunneling with NDR in
ZnO/Zn$_{0.8}$Mg$_{0.2}$O devices at room temperature
(Krishnamoorthy, Iliadis, Inumpudi, Choopun, \textit{Solid-State Electron.}\
vol.~46, pp.~1633--1637, 2002). This is the experimental grounding that
makes ZnO/Mg$_{0.3}$Zn$_{0.7}$O the natural primary stack for the SISPAD
paper~\cite{krishnamoorthy2002}, replacing our earlier
In$_2$O$_3$/$\kappa$-Ga$_2$O$_3$ candidate (see \S\ref{sec:stack}).

\textbf{Note:} An earlier version of this document cited ``Tampo et al.,
IEEE~NANO 2019'' for ZnO RTD NDR. That reference had a fabricated author
list (the DOI 10.1109/NANO46743.2019.8874570 exists but with different
authors). Krishnamoorthy~2002 is the verified earliest ZnO/MgZnO RTD NDR
demonstration.

\subsection{Datta, \textit{Quantum Transport: Atom to Transistor}}
Source for the analytic transverse-integration trick (``Fix-$A$'') that
turns the quasi-3D transverse degree of freedom into a prefactor in the
current integral, avoiding any artificial discrete-mode-basis
quantisation box (\S\ref{sec:method-fixA}).

% =====================================================================
\section{Theoretical Framework: NEGF with SCBA Inelastic Transport}
% =====================================================================

\subsection{Starting point: Meir--Wingreen current}
For a region coupled to two leads with retarded Green's function $G^R(E)$
and contact broadenings $\Gamma_{L,R}(E)$, the right-contact current in
the steady state is
\begin{equation}
I_R = \frac{e}{h}\int dE\,\mathrm{Tr}\bigl[\Sigma^{<}_R(E)\,G^{>}(E) -
\Sigma^{>}_R(E)\,G^{<}(E)\bigr].
\end{equation}
The lesser/greater Green's functions satisfy the Keldysh equation
\begin{equation}
G^{<,>}(E) = G^R(E)\,\bigl[\Sigma^{<,>}_L + \Sigma^{<,>}_R +
\Sigma^{<,>}_{\mathrm{ph}}\bigr]\,G^A(E),
\label{eq:keldysh}
\end{equation}
with $G^A = (G^R)^\dagger$ and contact distributions set by
$\Sigma^{<}_{L,R}(E) = i f_{L,R}(E)\,\Gamma_{L,R}(E)$ under Patil's
symmetric-bias convention $\mu_{L,R}(V) = E_F \pm V/2$.

\subsection{Self-consistent Born approximation (SCBA)}
For an electron-vibron coupling Hamiltonian
$H_{\mathrm{e-ph}} = \sum_\nu D_0 (b_\nu + b_\nu^\dagger)\,\mathcal M_\nu$,
the SCBA phonon self-energies at leading order in $D_0^2$ are
\begin{align}
\Sigma^{<,>}_{\mathrm{ph}}(E) &= \sum_\nu D_0^2\,\bigl[
\, N_\nu\,\mathcal M_\nu G^{<,>}(E \mp \hbar\omega_\nu)\mathcal M_\nu \nonumber\\
&\quad + (N_\nu + 1)\,\mathcal M_\nu G^{<,>}(E \pm \hbar\omega_\nu)\mathcal M_\nu
\bigr],\\
\Sigma^{R}_{\mathrm{ph}}(E) &= -\tfrac{i}{2}\bigl[\Sigma^{>}_{\mathrm{ph}}(E)
- \Sigma^{<}_{\mathrm{ph}}(E)\bigr]\;\text{(causal part only)},
\end{align}
with $N_\nu$ the Bose occupation at the electronic temperature. The full
SCBA iterates $G^{R,<,>} \to \Sigma^{R,<,>}_{\mathrm{ph}} \to G^{R,<,>}$
to self-consistency. Two implementation points where we have explicitly
departed from earlier (buggy) versions of our own code:
\begin{enumerate}[leftmargin=1.3em,itemsep=0pt]
\item The coupling in the code is $D^2$, not $D$ (fix applied 2026-03;
documented in \code{docs/COUPLING\_NORMALIZATION\_FIX.md}).
\item The retarded phonon self-energy must keep only the causal part
$-\tfrac{i}{2}(\Sigma^{>}-\Sigma^{<})$; an earlier fix corrected a sign
that was leaking anti-causal weight (\code{docs/BUG\_FIXES\_SUMMARY.md}
Fix~4).
\item Contact in-scattering $f_L\Gamma_L + f_R\Gamma_R$ must be added to
$\Sigma^{<}$ for inelastic modes (Fix~3); omitting it was the root cause
of the ``all molecules identical'' artefact in our earlier runs.
\end{enumerate}

\subsection{First Born vs SCBA --- what ``FBA'' means in our code}
``FBA'' in this codebase denotes a run with \code{max\_iter = 1}: the
phonon self-energy is computed once from a ballistic $G$, and the Green's
function is then \emph{recomputed once} before the current integral
(this ``final pass'' was the Fix~1 of early March that corrected a silent
identity of all FBA molecular currents). ``SCBA'' denotes full
self-consistency via Anderson mixing (DIIS). On the converged device the
two currents differ by $<$2\,\% for $D_0^2 = 0.1$~eV$^2$. Early in
the project FBA was the production solver because SCBA did not converge
near NDR; after implementing Anderson mixing (2026-04-14), SCBA became
the production solver for all SISPAD results. FBA remains available for
quick exploratory sweeps.

\subsection{Semi-infinite contact self-energies}
Each tight-binding contact contributes $\Sigma_{L,R}(E) = -t_0\,e^{\pm i k a}$
at the boundary site, with $\cos ka = 1 - (E - U_{L,R})/(2t_0)$. These are
the exact Bloch-continuation self-energies for a cosine dispersion and are
used throughout.

% =====================================================================
\section{The Rank-1 Projected SCBA (methodological core)}
\label{sec:rank1}
% =====================================================================

\subsection{Setup: quasi-3D RTD with a localized molecular vibron}
The elastic Hamiltonian separates,
$H_{\mathrm{el}} = H_z(z) + H_\perp(\mathbf r_\perp)$,
with a transverse box of area $A = L_y L_z$. The longitudinal direction is
discretised on $N_p$ tight-binding sites with spacing $a \sim$ 0.15--0.2~nm.
A single molecular vibron of frequency $\omega$ sits at
$\mathbf r_{\mathrm{mol}} = (z_{\mathrm{mol}}, \mathbf r_{\perp,
\mathrm{mol}})$ and couples through a local Fr\"ohlich-type coupling
$H_{\mathrm{e-ph}} = D_0 (b + b^\dagger) \!\int\! d^3 r\,
g(\mathbf r - \mathbf r_{\mathrm{mol}}) c^\dagger(\mathbf r) c(\mathbf r)$
with an isotropic Gaussian form factor of a \emph{single} length scale
$\sigma_{\mathrm{mol}} = 3$~\AA.

\subsection{Rank-1 theorem for the coupling matrix}
Projecting the vibron coupling onto the transverse eigenbasis at site $z$
gives
\begin{equation}
M_{nm,n'm'}(z) = D_0\,\chi(z - z_{\mathrm{mol}})\!\int\!d^2 r_\perp\,
\psi^*_{nm}(\mathbf r_\perp) g_\perp(\mathbf r_\perp - \mathbf r_{\perp,\mathrm{mol}}) \psi_{n'm'}(\mathbf r_\perp),
\end{equation}
with $\chi(z) \propto e^{-z^2/2\sigma_{\mathrm{mol}}^2}$. A Taylor expansion
of $\psi_{n'm'}$ around $\mathbf r_{\perp,\mathrm{mol}}$ combined with the
vanishing of odd Gaussian moments gives
\begin{equation}
\boxed{\;M_{nm,n'm'}(z) = D_0\,\chi(z - z_{\mathrm{mol}})\,u_{nm}\,u^*_{n'm'}\,\bigl[1 + O((\sigma_{\mathrm{mol}} k_{\perp,\max})^2)\bigr]\;}
\label{eq:rank1-factor}
\end{equation}
with $u_{nm} \equiv \psi_{nm}(\mathbf r_{\perp,\mathrm{mol}})$. For the
relevant geometry ($\sigma_{\mathrm{mol}} = 3$~\AA, $L_\perp = 1\,\mu$m,
$n_{\max} = 3$) the correction is $\sim 8 \times 10^{-6}$.

\textbf{Consequence.} The coupling operator is
$\mathcal M(z) = D_0 \chi(z)\,|u\rangle\langle u|$: \emph{rank-1 in mode
space to machine precision for realistic molecular sizes}. Crucially, the
rank-1 factorisation holds for \emph{any} transverse grid size, not only
the $3\times 3$ case used in the toy test. As $L_\perp$ increases or
$n_{\max} \to \infty$, the correction remains
$O((\sigma_{\mathrm{mol}} k_{\perp,\max})^2)$; in the continuum limit
($\Delta\varepsilon \ll k_BT$) the coupling is still rank-1 because
$\sigma_{\mathrm{mol}} \ll L_\perp$. The
$\sigma_{\mathrm{mol}} \to 0$ limit is not an idealisation but an
asymptotic result whose error scales as the sixth power of
$\sigma_{\mathrm{mol}}/(\text{transverse wavelength})$.

\subsection{Projected 1D Dyson equation}
Because $\mathcal M$ is rank-1 in modes, the SCBA self-energy inherits the
same rank-1 structure,
$\Sigma^{R,<,>}_{\mathrm{ph}}(z,z';E) = |u\rangle\langle u|\,
\tilde\sigma^{R,<,>}(z,z';E)$. Introducing the projected Green's function
$\tilde G^R = \langle u | G^R | u\rangle$, the full phonon-dressed Dyson
equation collapses to a single scalar-in-mode-space, 1D-in-longitudinal-
space equation for $\tilde G^R$. One SCBA iteration is a single 1D NEGF
solve and is $\sim\!4\times$ faster than the legacy ``4-inelastic +
5-coherent'' hybrid solver, while including all inter-mode off-diagonals
exactly.

\subsection{Exact coherent/inelastic current decomposition}
The Meir--Wingreen current decomposes rigorously (not heuristically) into
\begin{equation}
I = I_{\mathrm{coh}} + I_{\mathrm{inel}},
\end{equation}
with the coherent piece exactly the Datta quasi-1D Landauer sum (each mode
carries its bare transmission with its own Fermi shift by $\epsilon_{nm}$),
and the inelastic piece a bilinear in $\tilde\sigma^{<,>},\tilde G^R,
\tilde G^A$ weighted by $|u_{nm}|^2$. Because the split is a consequence of
the rank-1 factorisation of $\mathcal M$, it is a theorem rather than a
tunable partition.

\subsection{Exact parity theorem (symmetric-pixel limit)}
At a symmetric pixel with the molecule at its centre,
$u_{nm} = \psi_{nm}(L_y/2, L_z/2)$ vanishes identically whenever $n$ or $m$
is even. In a $3\times 3$ mode grid this decouples the 9 modes as
\emph{exactly four (odd,odd) inelastic} and \emph{exactly five elastic};
the earlier ``4+5 hybrid'' partition, previously justified by a
coupling-overlap threshold, is now a theorem. The continuum Datta
version of the argument gives an analogous parity selection rule on the
integral that survives the Fix-$A$ replacement.

\subsection{Multiple phonon modes: bulk + molecular}
\label{sec:bulk-mol}
In physical devices two classes of phonon participate simultaneously:

\textbf{(a) Bulk LO phonon.} The ZnO lattice phonon ($\hbar\omega_{\mathrm{LO}} = 72$~meV)
couples uniformly across the transverse plane. Its coupling matrix is
$\mathcal M^{\mathrm{bulk}}(z) = D_{\mathrm{bulk}}\,\chi_{\mathrm{bulk}}(z)\,\mathbf{I}_\perp$,
where $\mathbf{I}_\perp$ is the identity in mode space and $\chi_{\mathrm{bulk}} \equiv 1$
in the active region (barriers + well) and zero in the contacts. This is
\emph{diagonal} in mode space---no inter-mode scattering:
\begin{equation}
\Sigma^{\mathrm{bulk}}_{nm,n'm'} = \delta_{nm,n'm'}\,\tilde\sigma^{\mathrm{bulk}}_{nm}.
\end{equation}

\textbf{(b) Molecular vibron.} As derived in \S\ref{sec:rank1}, the localized
vibron coupling is rank-1: $\mathcal M^{\mathrm{mol}}(z) = D_0\,\chi(z)\,|u\rangle\langle u|$,
with inter-mode scattering captured exactly.

\textbf{Combined self-energy.} The total is $\Sigma = \Sigma^{\mathrm{bulk}} + \Sigma^{\mathrm{mol}}$,
i.e.\ diagonal + rank-1. The full 3D SCBA can in principle be solved in two stages:
(i)~9 independent bulk SCBAs for each transverse mode, then (ii)~one projected
molecular SCBA on $\tilde G = \langle u|G|u\rangle$ dressed by the bulk-converged
modes. In the \emph{continuum transverse limit} ($\Delta\varepsilon \sim 10^{-8}$~eV $\ll k_BT$
for the $10\,\mu$m pixel), all modes are degenerate and both self-energies reduce
to scalar 1D functions applied additively in a single SCBA loop:
\begin{equation}
\tilde\sigma^R(E) = \tilde\sigma^R_{\mathrm{bulk}}(E) + \tilde\sigma^R_{\mathrm{mol}}(E).
\end{equation}
This is the form implemented in \code{core/scba\_rank1\_keldysh.py}.

\textbf{Multiple molecular vibrons.} If a molecule has $K$ modes ($\omega_1, \ldots, \omega_K$),
each centred on the same spatial location, they share the same
$|u\rangle$ vector. The molecular self-energy becomes a sum of $K$ rank-1
terms, all with the same projector---hence still rank-1 as a matrix in
mode space. The implementation loops over phonon modes and accumulates
$\Sigma_{\mathrm{in/out}}$; the algebra is unchanged.

\textbf{Bulk coupling calibration.} Patil's reference code has $D_{\mathrm{bulk}} = 0$
(no bulk phonon). For ZnO we include the LO phonon with
$D^2_{\mathrm{bulk}} = 0.001$~eV$^2$ per site. The exact per-site
Fr\"ohlich coupling requires DFT calibration; 0.001 is defensible as
moderate coupling ($0.001 \times 35$ active sites $= 0.035$~eV$^2$ total).
Previous value 0.005 caused FBA near-singular spikes at 30\% of bias
points where the phonon-shifted resonance replica pushed $G^R$ poles onto
the real axis. At $D^2_{\mathrm{bulk}} = 0.001$, conservation residual is
$< 0.2\%$ at all bias points under SCBA.
A sensitivity sweep on $D^2_{\mathrm{bulk}}$ from 0.001 to 0.1~eV$^2$ showed
conservation degrading rapidly---the self-energy becomes comparable to the
bare Hamiltonian and the SCBA (or FBA) becomes numerically unstable.

\subsection{Datta ``Fix-$A$'' transverse integration}
\label{sec:method-fixA}
For a separable parabolic-subband quasi-3D Hamiltonian the transverse
degree of freedom integrates analytically (Datta 2005, Eq.~9.3.14):
\begin{equation}
I = \frac{2e}{h}\,A\,\frac{m^* k_B T}{2\pi\hbar^2}
\!\int\! dE_z\, T(E_z)\,\Bigl\{\log\bigl[1\!+\!e^{(E_F - E_z + V/2)/k_BT}\bigr]
- \log[1\!+\!e^{(E_F - E_z - V/2)/k_BT}]\Bigr\},
\end{equation}
so that the pixel area $A$ enters only as a multiplicative prefactor, not
as a quantisation box. For the 10~\si{\um}$\times$10~\si{\um} pixel chosen
here (see \S\ref{sec:stack}) the transverse subband spacing is $\sim 40$~neV,
far below $k_BT$, and the continuum integral is exact. The discrete mode
grid that predated this fix was $3\!\times\!3$ and has been abandoned for
the production runs.

\subsection{Closed-form analytic $d^2I/dV^2$}
Differentiating $I_R(V) = (e/h)\int dE\,\mathrm{Tr}[\cdots]$ twice in $V$
under the symmetric bias convention and using the optical-theorem identity
$\Gamma_L + \Gamma_R = i(\Sigma^R - \Sigma^A)$, which at $D_0 \to 0$ becomes
$T_{RA} = T_{RL} + T_{RR}$, gives
\begin{equation}
\boxed{\;\frac{d^2 I_R}{dV^2}(V) = \frac{e^2}{h}\cdot\frac{1}{4}\!\int dE\,
\bigl[f_L''(E,V)\,T_{RL}(E) + f_R''(E,V)\,(T_{RR}(E) - T_{RA}(E))\bigr]\;}
\label{eq:d2i}
\end{equation}
where $f''_{L,R}$ are the second derivatives of the Fermi functions
evaluated at $\mu_{L,R}(V)$, and the kernels $T_{RL}, T_{RR}, T_{RA}$ are
built once from the converged $(G^R, \Gamma_L, \Gamma_R)$ at each bias.
The crucial property is that this formula eliminates all numerical
$V$-grid double-differentiation, which was the dominant source of noise
in our earlier IETS plots. At $D_0 = 0$ it reduces exactly to the
Landauer coherent second derivative; this reduction is unit-tested to
$10^{-6}$ (\code{tests/test\_analytic\_d2idv2.py}).

% =====================================================================
\section{Material Stack Selection: BEOL-Compatible Oxide RTDs}
\label{sec:stack}
% =====================================================================

\subsection{Why low-barrier oxides}
Traditional oxide dielectrics (Al$_2$O$_3$, HfO$_2$, SiO$_2$) have
conduction-band offsets of $\sim\!2$--3~eV, putting any RTD NDR feature at
4--5~V bias where \emph{barrier-top} leakage dominates and all molecular
IETS information is washed out. Patil's operating window is 0.3--0.8~V,
corresponding to molecular vibrational energies of 90--800~meV. To
replicate that window in a BEOL-compatible oxide we need
$\Delta E_c \sim 0.5$~eV. The systematic literature search in
\code{docs/LOW\_BARRIER\_OXIDE\_IETS.md} identified four viable oxide pairs:
\begin{center}
\small
\begin{tabular}{lccl}
\toprule
System & $\Delta E_c$ (eV) & ALD $T$ & Notes\\
\midrule
In$_2$O$_3$/$\beta$-Ga$_2$O$_3$      & 0.07 & $<$400\,\si{\celsius} & Too low, weak NDR in sim\\
In$_2$O$_3$/$\kappa$-Ga$_2$O$_3$     & 0.45 & $<$400\,\si{\celsius} & Good sim NDR, metastable $\kappa$\\
ZnO/Mg$_{0.3}$Zn$_{0.7}$O             & 0.47 & $<$350\,\si{\celsius} & Good sim NDR, experimentally demoed\\
ZnO/$\alpha$-Ga$_2$O$_3$              & 0.61 & $<$400\,\si{\celsius} & Higher CBO, less attractive\\
\bottomrule
\end{tabular}
\end{center}

\subsection{Primary stack decision: ZnO/Mg\textsubscript{0.3}Zn\textsubscript{0.7}O}
The primary stack for the SISPAD submission was swapped from
In$_2$O$_3$/$\kappa$-Ga$_2$O$_3$ to ZnO/Mg$_{0.3}$Zn$_{0.7}$O on 2026-04-07
(\code{docs/STACK\_DECISION.md}). The rationale:
\begin{itemize}[leftmargin=1.3em,itemsep=0pt]
\item Every physics number is within a few percent of the In$_2$O$_3$
stack: $\Delta E_c = 0.47$ vs $0.45$~eV, $m^* = 0.28\,m_0$ vs $0.30\,m_0$,
$\hbar\omega_{\mathrm{LO}} \approx 72$ vs $70$~meV.
\item Every \emph{reviewer-facing} axis favours MgZnO: stable wurtzite
vs metastable $\kappa$; multi-source CBO vs single XPS source; decades-old
ALD process vs an emerging one; \emph{experimentally demonstrated RTD with
NDR}~\cite{krishnamoorthy2002} vs none in the $\kappa$-Ga$_2$O$_3$ system.
\item The CBO is tunable via $x_{\mathrm{Mg}}$, scaling as $1.57\,x$; this
adds an experimental degree of freedom for sensitivity tuning.
\end{itemize}
The In$_2$O$_3$/$\kappa$-Ga$_2$O$_3$ stack becomes the Section III
comparison data: a side-by-side FBA run on both stacks confirms that the
methodology result transfers between oxide pairs.

\subsection{Layer stack (primary device)}
\begin{center}
\small
\begin{tabular}{cll}
\toprule
Layer & Material & Thickness\\
\midrule
emitter contact & $n^+$-ZnO, $N_D = 10^{25}\,\si{\per\meter\cubed}$ & 10\,nm\\
emitter barrier & Mg$_{0.3}$Zn$_{0.7}$O & 3.0\,nm\\
well            & ZnO                   & 3.0\,nm\\
collector barrier & Mg$_{0.3}$Zn$_{0.7}$O & 1.5\,nm\\
collector contact & $n^+$-ZnO, $N_D = 10^{25}\,\si{\per\meter\cubed}$ & 10\,nm\\
\bottomrule
\end{tabular}
\end{center}
Transverse area $A = 10\,\si{\um} \times 10\,\si{\um} = 10^{-7}\,\si{\cm\squared}$,
matching typical e-nose sensor-pixel footprints and published oxide RTD
mesas. $A$ enters only as a Datta prefactor (\S\ref{sec:method-fixA}).

\subsection{Molecule placement and the Einstein vibron assumption}
For the primary runs the molecular form factor is a Gaussian of width
$\sigma_{\mathrm{mol}} = 3$~\AA\ centered on the emitter-barrier midpoint
(the point of maximum wavefunction overlap with the resonant level). The
molecular vibron is treated as an Einstein oscillator with a single mode
energy $\hbar\omega$, Bose-distributed at the electronic temperature, with
no intrinsic linewidth. Physical lifetime broadening is a known omission
and is on the follow-up list (\S\ref{sec:roadmap}).

% =====================================================================
\section{Simulation Methodology}
\label{sec:methodology}
% =====================================================================

\subsection{Numerical grids}
\begin{itemize}[leftmargin=1.3em,itemsep=0pt]
\item Longitudinal: tight-binding grid with $a = 0.20$~nm,
$N_p = 135$ sites for the symmetric 2/3/2~nm primary device. Hopping
$t_0 = \hbar^2/(2m^* a^2) = 3.402$~eV.
\item Energy: $dE = 2$~meV over $[-0.25, +0.50]$~eV (376 points) for
production runs; $dE = 0.5$~meV for benchmark.
\item Bias: 201 points over 0--800~mV for SISPAD primary (4~mV spacing).
\item Broadening: $\eta = 10^{-12}$~eV (effectively zero; the semi-infinite
contact self-energies provide all the necessary broadening). Early runs
used $\eta = 20$~meV, which artificially broadened the resonance and was
removed.
\item Solver: SCBA with Anderson mixing (DIIS), history depth $M=8$,
damping $\beta = 0.3$, tolerance $10^{-4}$. See
\S\ref{sec:fba-vs-scba} for details.
\item Temperature: 300~K production; additional 10, 77, 150~K runs for
the temperature sensitivity panel.
\end{itemize}

\subsection{SCBA convergence: Anderson mixing (DIIS)}
\label{sec:fba-vs-scba}
The self-consistent Born approximation (SCBA) iterates
$G \to \Sigma_{\mathrm{ph}} \to G$ to self-consistency.
Early attempts with simple linear mixing ($\Sigma^{(k+1)} = (1-\alpha)\Sigma^{(k)} + \alpha\Sigma^{\mathrm{new}}$,
$\alpha = 0.3$--$0.5$) failed to converge near NDR:
the sharp resonance pole sweeps through the integration window under small
$\Sigma$ perturbations, causing the iteration to oscillate.

\textbf{Anderson mixing (DIIS)}~\cite{anderson1965} solved this completely
(implemented 2026-04-14 in \code{core/scba\_rank1\_keldysh.py}). The algorithm stores
the last $M = 8$ input/residual vector pairs, where each vector is the
packed diagonal of $\Sigma^{<,>}_{\mathrm{ph}}$ (a $4 N_E N_p$-dimensional
real vector). At each iteration, a constrained least-squares problem
\begin{equation}
\min_{\{\alpha_i\}} \Bigl\|\sum_i \alpha_i r_i\Bigr\|^2 \quad \text{s.t.} \quad \sum_i \alpha_i = 1
\end{equation}
is solved via a Lagrange multiplier on the $(M\!+\!1) \times (M\!+\!1)$
augmented system. The next iterate is
$x_{k+1} = \sum_i \alpha_i (x_i + \beta\, r_i)$ with $\beta = 0.3$.
Tikhonov regularisation ($10^{-12} I$) stabilises the overlap matrix.

\textbf{Results:} All bias points (0--800~mV) now converge in 10--60
iterations with current conservation $|I_L + I_R|/\max(|I|) < 0.2\%$,
including the previously fatal NDR region ($V = 0.49$--$0.6$~V). The
converged I(V) is inherently smooth, enabling raw numerical
$d^2I/dV^2 = \mathrm{diff}(\mathrm{diff}(I))/dV^2$ with no post-processing---exactly
as in Patil's MATLAB code (lines~148--149 of \code{rtd2modes\_1d.m}).

\textbf{Singular-matrix regularisation.} At rare $(E,V)$ points where the
phonon self-energy pushes $G^R$ poles near the real axis,
\code{np.linalg.solve} can still fail with \code{LinAlgError}. A
\code{\_robust\_solve()} wrapper adds $10^{-8}j \cdot I$ regularisation as
a fallback, without perceptibly changing the current elsewhere.

\textbf{FBA remains available} as \code{max\_iter = 1} for quick
exploratory sweeps; it includes a ``FINAL PASS'' that recomputes $G$ with
the phonon $\Sigma$ before the current integral (without this, all
molecules produce identical ballistic currents). However, SCBA with
Anderson mixing is now the production solver for all SISPAD results.

\subsection{Performance engineering (important)}
A $\sim$100$\times$ speedup on the rank-1 inner loop came from two fixes:
\begin{enumerate}[leftmargin=1.3em,itemsep=0pt]
\item Vectorising the per-energy Python loop in the $G^R$ construction
into a batched \code{np.linalg.solve} on arrays of shape
$(N_E, N_p, N_p)$.
\item Replacing \code{np.einsum('kij,kjm,kml->kil', ...)} in the
$G^{<,>}$ computation by an explicit
\code{np.matmul(np.matmul(G\_R, sig), G\_A)}. Under NumPy~2.4 / Python~3.14
the 3-way einsum path was falling back to an unoptimised kernel;
\code{matmul} delegates directly to BLAS.
\end{enumerate}
After these fixes a 201-bias rank-1 SCBA sweep at production settings
(10--60 Anderson-mixing iterations per bias point) completes in
$\sim$3--5~h on a single workstation core; a single FBA sweep
(\code{max\_iter = 1}) takes $\sim$1.5~h.

\subsection{Code map}
\small
\begin{tabular}{ll}
\toprule
Path & Role\\
\midrule
\code{core/scba\_solver\_hybrid.py}  & Legacy 4-inelastic + 5-coherent hybrid solver\\
\code{core/scba\_solver\_rank1.py}   & Rank-1 projected SCBA driver\\
\code{core/scba\_rank1\_keldysh.py}  & Single-bias rank-1 Keldysh wrapper (batched)\\
\code{core/self\_energy.py}          & Phonon self-energies ($D^2$ convention; causal $\Sigma^R$)\\
\code{core/iets\_analytic.py}        & Closed-form analytic $d^2I/dV^2$ (Eq.~\ref{eq:d2i})\\
\code{config/device\_library.py}     & Device layer stacks and material parameters\\
\code{config/molecular\_database.py} & 20 molecules (10 synthetic + 10 DFT-derived VOCs)\\
\code{run/run\_rank1\_sweep.py}      & Dense-$V$ production driver (SISPAD primary)\\
\code{run/run\_iets.py}              & Legacy production driver (hybrid solver)\\
\code{run/run\_oxide\_baseline.py}   & Baseline elastic I-V sweeps for stack selection\\
\code{papers/rtd2modes\_1d.m}        & Patil 1D reference MATLAB code\\
\code{tests/test\_patil\_reference.py} & Patil 1D benchmark harness\\
\code{tests/test\_analytic\_d2idv2.py} & Landauer-limit $10^{-6}$ unit test\\
\code{tests/test\_rank1\_vs\_full.py}  & Rank-1 vs brute-force coupled-SCBA toy test\\
\bottomrule
\end{tabular}

% =====================================================================
\section{Validation}
\label{sec:validation}
% =====================================================================

\subsection{Tier-1 checklist (SISPAD gating items)}
Four technical items were declared gating before the SISPAD submission
could be defended in front of reviewers
(\code{docs/SISPAD\_CHECKLIST.md}). All four closed locally on 2026-04-09:
\begin{center}
\small
\begin{tabular}{llll}
\toprule
Tier-1 item & Description & Status & Code\\
\midrule
$\alpha$ & Closed-form analytic $d^2I/dV^2$ & \textbf{closed} & \code{core/iets\_analytic.py}\\
$\beta$  & FBA vs SCBA separation (naming, numerics) & \textbf{closed} & \code{scba\_solver\_hybrid.py}\\
$\gamma$ & Patil 1D benchmark $\le 1\%$ at NDR peak & \textbf{closed (0.55\%)} & \code{test\_patil\_reference.py}\\
$\delta$ & Per-bias current conservation reported   & \textbf{closed} & \code{run\_rank1\_sweep.py}\\
\bottomrule
\end{tabular}
\end{center}
Two tightening follow-ups are tracked but do not block submission:
(1)~a residual $\sim 9\%$ gap between Python and Octave at a Patil
subresonance (NDR peak matches to 0.55\% regardless), and (2)~a median
$\sim\!20\%$ current-conservation residual at the benchmark that wants
\code{max\_iter $\to$ 500} at lower mix.

\subsection{Patil 1D benchmark}
The rank-1 Keldysh wrapper reproduces the Patil reference $I_R(V)$ NDR
peak to 0.55\%:
$I_R = 5.667 \times 10^{-10}$~A (rank-1) vs $5.698 \times 10^{-10}$~A
(Patil MATLAB) at $V = 0.336$~V. All mesh, bias, contact and phonon
parameters are held identical. Both contacts are exposed so
$\delta(V) = |I_L + I_R|/\max(|I_L|,|I_R|)$ is read directly from the
Meir--Wingreen trace at each bias. See
\code{tests/patil\_reference\_1d.py} and
\code{tests/patil\_reference\_1d.npz}.

\subsection{Optical-theorem unit test}
At $D_0 = 0$ the analytic $d^2I/dV^2$ formula~(\ref{eq:d2i}) must reduce
to the Landauer coherent second derivative by the optical theorem
$T_{RA} = T_{RL} + T_{RR}$. This is unit-tested to
$|\Delta|/\max < 10^{-6}$ in \code{tests/test\_analytic\_d2idv2.py}.

\subsection{Rank-1 vs brute-force coupled-SCBA toy test}
\code{tests/test\_rank1\_vs\_full.py} runs a small 2-mode $\times$ 3-site
toy problem with the same parameters through (a) a brute-force coupled-SCBA
with the full $\mathcal M$ matrix and (b) the rank-1 projected solver, and
checks they agree to machine precision. This is the unit confirmation of
Eq.~(\ref{eq:rank1-factor}).

% =====================================================================
\section{Results}
% =====================================================================

\subsection{Comparison-stack FBA sweep (In\texorpdfstring{$_2$}{2}O\texorpdfstring{$_3$}{3}/\texorpdfstring{$\kappa$}{kappa}-Ga\texorpdfstring{$_2$}{2}O\texorpdfstring{$_3$}{3})}
\label{sec:results-comparison}
Baseline, Mol\_A, Mol\_B and Mol\_AB were run on the In$_2$O$_3$/$\kappa$-Ga$_2$O$_3$
symmetric stack with the legacy hybrid FBA solver on an 81-point
$V$-grid from 0 to 400~mV at 300~K (cloud runs, 2026-03-25--03-29).
Peak-signal summary:
\begin{center}
\small
\begin{tabular}{lccc}
\toprule
Species & $I_{\mathrm{peak}}$ (nA) & $|d^2I/dV^2|_{\mathrm{peak}}$ ($\mu$A/V$^2$) & $|\Delta D|_{\mathrm{peak}}$ ($\mu$A/V$^2$)\\
\midrule
Baseline & 49.64 & 465.02 & ---  \\
Mol\_A   & 49.80 & 465.10 & 1.38 \\
Mol\_B   & 49.67 & 465.02 & 1.05 \\
Mol\_AB  & 49.88 & 465.10 & 2.43 \\
\bottomrule
\end{tabular}
\end{center}
Two observations are worth stressing: the absolute currents are within
$<$0.5\% of each other across species (the molecule perturbs the
resonance but does not destroy it), and the additivity check
$|\Delta D|_{AB} \approx |\Delta D|_A + |\Delta D|_B$ is exact within
numerical noise. The latter is a direct demonstration of the analytic
second-derivative formula~(\ref{eq:d2i}) acting linearly on a mixed
analyte. Plots: \code{results/2026-04-09/plots/cloud\_primary\_*.png}.

\subsection{Primary-stack rank-1 SCBA sweep (ZnO/MgZnO, symmetric 2/3/2~nm)}
\label{sec:results-primary}
The production sweep uses the rank-1 Keldysh SCBA with Anderson mixing on the
ZnO/Mg$_{0.3}$Zn$_{0.7}$O symmetric 2/3/2~nm device: 201 bias points from
0 to 800~mV (4~mV spacing), $dE = 2$~meV, $\eta = 10^{-12}$~eV, $E_F = 20$~meV, 300~K.
Bulk LO phonon ($\hbar\omega = 72$~meV, $D^2_{\mathrm{bulk}} = 0.001$~eV$^2$/site,
$\chi \equiv 1$ in active region) is always present; molecular modes are
added with $D^2_0 = 0.1$~eV$^2$ and Gaussian $\chi(z)$ at the emitter
barrier.

Key results:
\begin{itemize}[leftmargin=1.3em,itemsep=0pt]
\item \textbf{NDR at $V_{\mathrm{res}} \approx 560$~mV} with peak current
$\sim 69$~nA (Baseline) and peak-to-valley ratio $\sim 6$. The second
quantised state lies above the CBO, confirming single-resonance physics.
\item \textbf{Phonon sideband at $V \approx 2 \hbar\omega_{\mathrm{LO}}/e = 144$~mV}:
a dip in the Baseline I--V where phonon emission opens. This is real physics,
not numerical noise.
\item \textbf{Molecular discrimination:} Mol\_A (100~meV) and Mol\_B (180~meV)
produce distinct features in $d^2I/dV^2$ in the 200--500~mV range.
Mol\_AB (both modes) shows linear superposition:
$\Delta D_{AB} \approx \Delta D_A + \Delta D_B$.
\item \textbf{Current conservation:} $|I_L + I_R|/\max|I| < 0.2$\% at all
bias points under SCBA with Anderson mixing, including the NDR region.
\item \textbf{Smoothness:} The self-consistent SCBA produces inherently
smooth $I(V)$, so the raw numerical second derivative
$d^2I/dV^2 = \mathrm{diff}(\mathrm{diff}(I))/dV^2$ is directly usable as a
molecular fingerprint---exactly as in Patil's MATLAB code (lines~148--149 of
\code{rtd2modes\_1d.m}), with no smoothing or post-processing.
\item \textbf{$d^2I/dV^2$ fingerprinting:} Numerical double-differentiation
of the converged SCBA current reveals molecule-specific peak positions
corresponding to vibrational mode energies. The discrimination metric
$\Delta D(V) = d^2I_{\mathrm{mol}}/dV^2 - d^2I_{\mathrm{BL}}/dV^2$
isolates molecular signatures by subtracting the bulk-phonon background.
Mol\_AB exhibits approximate superposition of the individual Mol\_A and
Mol\_B signatures---a direct consequence of the rank-1 coupling structure
at weak coupling ($D_0^2 = 0.1$~eV$^2$).
\end{itemize}

\subsection{Temperature dependence (Mol\_A on symmetric 2/3/2~nm)}
\label{sec:results-temp}
I--V and $d^2I/dV^2$ sweeps for Mol\_A at 10~K, 77~K, 150~K and 300~K on the same
device. Key observations:
\begin{itemize}[leftmargin=1.3em,itemsep=0pt]
\item \textbf{I--V sharpening:} The NDR peak sharpens and increases at lower
$T$: $\sim$110~nA at 10~K vs $\sim$69~nA at 300~K, due to reduced thermal
smearing of the Fermi functions. PVR improves from $\sim$6 at 300~K to
$\sim$15 at 10~K.
\item \textbf{IETS feature sharpening:} The $d^2I/dV^2$ features
sharpen monotonically as $T$ decreases, confirming that the RTD resonance
provides an energy filter tighter than $k_BT$: even at 300~K the vibrational
peak is resolved, though thermally broadened. The 10~K peak-to-peak
amplitude is $\sim 5\times$ larger than at 300~K.
\item \textbf{Room-temperature viability:} The central claim of the project
--- that RTD-IETS can work at 300~K --- is directly supported: the Mol\_A
100~meV vibrational signature remains clearly resolvable in both $\Delta I$
and $\Delta D(V)$ at room temperature. The RTD resonance acts as an energy
filter tighter than $k_BT = 26$~meV at 300~K, preserving vibrational
features that would be washed out in a conventional MIM IETS junction.
\end{itemize}

\subsection{Figures}
\begin{figure}[h]
\centering
\placeholder{0.72\textwidth}{5.2cm}
\caption{\textbf{Device and methodology at a glance.} Symmetric
ZnO/Mg$_{0.3}$Zn$_{0.7}$O RTD (2/3/2~nm), band profile at $V = 0$ and
$V = V_{\mathrm{res}}$, Gaussian $\chi(z)$ with $\sigma_{\mathrm{mol}} = 3$~\AA,
and the rank-1 projected SCBA loop. See \S\ref{sec:rank1}.}
\end{figure}
\begin{figure}[h]
\centering
\placeholder{0.72\textwidth}{5.0cm}
\caption{\textbf{Patil 1D benchmark.} Rank-1 Keldysh wrapper vs Patil's
reference MATLAB code \code{rtd2modes\_1d.m} at the NDR peak; agreement
0.55\%. Inset shows per-bias $|I_L+I_R|/\max|I|$. See \S\ref{sec:validation}.}
\end{figure}
\begin{figure}[h]
\centering
\placeholder{0.72\textwidth}{5.0cm}
\caption{\textbf{I--V characteristics} of the symmetric 2/3/2~nm
ZnO/MgZnO RTD at 300~K for Baseline, Mol\_A, Mol\_B and Mol\_AB
(rank-1 SCBA with Anderson mixing, 201 bias points, 0--800~mV). NDR at $\approx$560~mV.
See \S\ref{sec:results-primary}.}
\end{figure}
\begin{figure}[h]
\centering
\placeholder{0.72\textwidth}{5.0cm}
\caption{\textbf{Numerical $d^2I/dV^2$ and discrimination metric} $\Delta D(V)$
at 300~K (SCBA). $d^2I/dV^2$ computed by raw numerical double-differentiation of
the converged SCBA current; no smoothing. Mol\_A peaks near 100~mV, Mol\_B
near 180~mV. Mol\_AB approximately equals the linear sum.
See \S\ref{sec:results-primary}.}
\end{figure}
\begin{figure}[h]
\centering
\placeholder{0.72\textwidth}{5.0cm}
\caption{\textbf{Temperature dependence of $d^2I/dV^2$} for Mol\_A at 10, 77,
150 and 300~K. IETS features sharpen at lower $T$; room-temperature features
remain resolvable. See \S\ref{sec:results-temp}.}
\end{figure}

% =====================================================================
\section{Open Issues and Limitations}
% =====================================================================

\subsection{Physics approximations we knowingly make}
\begin{itemize}[leftmargin=1.3em,itemsep=0pt]
\item \textbf{Einstein vibron.} The molecular mode is treated as an infinitely
narrow oscillator. Physical lifetime broadening $\gamma_\nu$ adds an extra
linewidth to the IETS feature; this is quantifiable and on the roadmap.
\item \textbf{Linear bias drop.} $U_{\mathrm{bias}}(z)$ is currently linear
between contacts. For the resonance positions and shapes this is a
$\sim$5--10\% effect; a self-consistent Poisson loop is planned.
\item \textbf{Image-charge renormalisation} of the vibron frequency by the
oxide environment ($\sim 50$~meV shift) is noted but not corrected, because
our SISPAD analytes are synthetic.
\item \textbf{Point-molecule placement.} We report the single-molecule
optimal-placement response. The ensemble-averaged response (random
position over the pixel) is $\sim\!\tfrac12$ of the optimal value by
rough estimate and will populate the full paper's appendix.
\item \textbf{Bulk phonon coupling calibration.} We include the ZnO LO phonon
with $D^2_{\mathrm{bulk}} = 0.001$~eV$^2$/site (reduced from an earlier
0.005 value that caused FBA near-singular spikes; see Fix~9/11 in Appendix~B).
A rigorous first-principles
value requires computing the Fr\"ohlich coupling per tight-binding site
from the ZnO macroscopic parameters ($\alpha_F \approx 1.06$,
$\epsilon_\infty = 3.72$, $\epsilon_s = 8.12$); this is on the follow-up list.
The Baseline trace serves as a sensitivity reference.
\item \textbf{No self-consistent Poisson (Hartree) potential.} The bias
profile $U_{\mathrm{bias}}(z)$ is currently a \emph{linear} drop from
$+V/2$ in the emitter contact to $-V/2$ in the collector, with no
self-consistent Hartree loop. In a physical RTD, charge accumulates in the
quantum well at resonance, creating a space-charge potential that:
\begin{enumerate}[leftmargin=1em]
\item Shifts the resonant level position (and hence $V_{\mathrm{NDR}}$)
by $\sim$5--15\% depending on doping and geometry;
\item Reshapes the potential from linear to staircase-like, with most of
the drop across the collector barrier;
\item Modifies the peak-to-valley current ratio.
\end{enumerate}
The self-consistent Poisson loop is the most important missing piece for
quantitative accuracy. The standard approach is: (i)~compute the electron
density $n(z)$ from $G^<$ at each NEGF iteration, (ii)~solve the 1D
Poisson equation $\nabla^2 \phi = -\rho/\epsilon$ with the doping profile
as background charge, (iii)~update $U_{\mathrm{bias}}(z) = -e\phi(z)$ and
repeat until the potential converges. Typical convergence requires 5--20
outer Poisson iterations with under-relaxation. This is planned as the
first extension for the IEEE TED paper.

For the SISPAD submission, the linear-drop approximation is acceptable
because: (a)~the symmetric 2/3/2~nm structure with equal doping
distributes the drop more evenly than an asymmetric design;
(b)~the IETS discrimination metric $\Delta D(V)$ is a differential
quantity that is less sensitive to the absolute potential profile than
the absolute current; (c)~the synthetic analytes allow relative comparisons
on the same device without requiring quantitative current-matching to
experiment.
\end{itemize}

\subsection{Numerical tightening items}
\begin{itemize}[leftmargin=1.3em,itemsep=0pt]
\item $\gamma$-tighten: $\sim 9$\% Python/Octave gap at a Patil subresonance
needs a root-cause investigation (contact-lead Bloch continuation sign,
most likely).
\item \textbf{SCBA convergence near NDR (resolved 2026-04-14).}
Each SCBA iteration requires a full $(N_E \times N_p \times N_p)$ matrix
solve ($376 \times 135 \times 135 \approx 7\times 10^6$ elements); one
iteration takes $\sim$0.7~s locally. Near the NDR voltage, the resonance
pole sweeps through the integration window under small $\Sigma$ perturbations,
causing simple linear mixing to oscillate rather than converge. The fix
progression: Fix~10 switched to a relative convergence criterion (60\% of
points converged); then Anderson mixing (DIIS) with history depth $M=8$
and damping $\beta = 0.3$ achieved convergence at \emph{all} bias points
in 10--60 iterations, with current conservation $< 0.2\%$.
\textbf{Resolution:} SCBA with Anderson mixing is now the production solver
for all SISPAD results (see \S\ref{sec:fba-vs-scba}).
\item Energy-grid aliasing: documented in
\code{docs/ENERGY\_GRID\_ALIASING.md}. Current $dE = 2$~meV is safe for
the synthetic analytes; realistic VOCs with $\hbar\omega < 50$~meV will
require finer sampling.
\end{itemize}

% =====================================================================
\section{Roadmap}
\label{sec:roadmap}
% =====================================================================

\subsection{SISPAD 2026 submission (deadline 2026-04-22)}
\begin{itemize}[leftmargin=1.3em,itemsep=0pt]
\item \textbf{[done]} SCBA with Anderson mixing production sweep on symmetric
2/3/2~nm primary stack (201~pts, 0--800~mV, 4~mV spacing, Baseline +
Mol\_A + Mol\_B + Mol\_AB at 300~K; Mol\_A at 10, 77, 150~K).
Data files in \code{results/sispad\_scba\_2026-04-14/}.
\item \textbf{[done]} Tier-1 $\alpha$, $\beta$, $\gamma$, $\delta$.
\item \textbf{[done]} SCBA convergence solved via Anderson mixing (DIIS).
All bias points converge in 10--60 iterations with conservation $< 0.2\%$
(\S\ref{sec:fba-vs-scba}).
\item \textbf{[done]} All 8 SISPAD figures generated from SCBA production data.
\item \textbf{[done]} Tampo~2019 reference replaced with verified
Krishnamoorthy~2002; fabricated references removed.
\item \textbf{[pending]} Complete author/affiliation block, acknowledgements
in \code{papers/IEEE-conference-template-062824/IEEE-conference-template-062824.tex}.
\item \textbf{[pending]} Regenerate figures from 201-pt production data when
runs complete (51-pt pilot data currently used).
\end{itemize}

\subsection{Full IEEE TED journal paper (target: 6--8 weeks post-SISPAD)}
Extensions beyond the SISPAD 2-page abstract:
\begin{enumerate}[leftmargin=1.3em,itemsep=2pt]
\item \textbf{Self-consistent Poisson loop} on $U_{\mathrm{bias}}(z)$.
\item \textbf{Realistic VOC analytes} from DFT-derived vibrational spectra:
benzene, thiophene, naphthalene, dimethyl sulfide, acetone, and four
others already in \code{config/molecular\_database.py}.
\item \textbf{Ensemble averaging} of molecular position over the pixel.
\item \textbf{Temperature study} from 10--300~K, showing the cross-over
from Patil's cryogenic regime to room temperature.
\item \textbf{Bulk-phonon sensitivity} (sweeping $D_{\mathrm{bulk}}$).
\item \textbf{Image-charge correction} via a simple image-plane dielectric
model.
\item \textbf{Noise floor} quantification using a realistic
$1/f$ contact fluctuation model and a lock-in detector of stated
time constant, so $\Delta D(V)$ can be reported as an SNR.
\item \textbf{Array-level discrimination.} A pixel-array pattern-recognition
stage (linear SVM or random forest) on a panel of analytes, to quantify
array-level classification accuracy at 300~K.
\item \textbf{Direct comparison with the hybrid 4+5 solver} for a handful
of cases, documenting the speedup and confirming the parity theorem
numerically on finite pixels.
\item \textbf{Full SCBA convergence --- done.} Anderson mixing (DIIS) was
implemented and deployed for the SISPAD submission (see \S\ref{sec:fba-vs-scba}).
All bias points converge in 10--60 iterations. The SCBA recovers the FBA
result in the $D_0 \to 0$ limit and shows $<$2\% deviation at
$D_0^2 = 0.1$~eV$^2$. This item is closed; the TED paper will present
SCBA as the production solver throughout.
\end{enumerate}

\subsection{Scientific extensions beyond the first journal paper}
\begin{enumerate}[leftmargin=1.3em,itemsep=4pt]
\item \textbf{Self-consistent Poisson loop on $U_{\mathrm{bias}}(z)$.}
The linear-drop approximation shifts $V_{\mathrm{NDR}}$ by $\sim$5--15\%.
The standard approach iterates: compute $n(z)$ from $G^<$, solve 1D
Poisson $\nabla^2\phi = -\rho/\epsilon$ with the doping background, update
$U_{\mathrm{bias}} = -e\phi(z)$, repeat until self-consistency (typically
5--20 outer iterations with under-relaxation). The SCBA phonon loop runs
\emph{inside} each Poisson iteration, producing a doubly self-consistent
simulation. This is the most important missing piece for quantitative
comparison with experiment. A scaffold exists in
\code{quantum\_enose\_poisson/} but is not yet integrated. For the SISPAD
submission, the linear-drop approximation is acceptable because:
(a)~the symmetric structure distributes the drop more evenly;
(b)~the discrimination metric $\Delta D(V)$ is differential and less
sensitive to the absolute profile; (c)~the synthetic analytes allow
relative comparisons without current-matching to experiment.

\item \textbf{DFT-calibrated Fr\"ohlich coupling.}
The phenomenological $D^2_{\mathrm{bulk}} = 0.001$~eV$^2$/site is an
order-of-magnitude estimate. A rigorous first-principles value requires
computing the Fr\"ohlich coupling per tight-binding site from the ZnO
macroscopic parameters: $\alpha_F \approx 1.06$,
$\epsilon_\infty = 3.72$, $\epsilon_s = 8.12$,
$\hbar\omega_{\mathrm{LO}} = 72$~meV. The per-site coupling is
$D^2 = (\alpha_F / N_{\mathrm{sites}}) \hbar\omega_{\mathrm{LO}} \times
(\text{geometry factor})$; the geometry factor depends on the
tight-binding wavefunction normalisation and the polar optical phonon
displacement field. A companion DFT calculation (or published
first-principles GW self-energy data for ZnO) would close this gap.
Similarly, the molecular coupling $D_0^2 = 0.1$~eV$^2$ is Patil's
phenomenological value; replacing it with DFT dipole-moment derivatives
for each analyte would close the loop between chemistry and transport.

\item \textbf{Einstein vibron lifetime broadening.}
The molecular vibron is currently an infinitely narrow oscillator. Real
molecules have vibrational lifetimes $\tau_\nu \sim 1$--10~ps
(corresponding to linewidths $\gamma_\nu \sim 0.1$--1~meV), set by
anharmonic phonon--phonon coupling and molecule--surface interactions. The
IETS peak width at low $T$ is dominated by $\gamma_\nu$ rather than
$k_BT$. Promoting the Einstein oscillator to a Lorentzian spectral
function $\delta(\omega - \omega_\nu) \to
\gamma_\nu/[\pi((\omega-\omega_\nu)^2 + \gamma_\nu^2)]$ in the phonon
propagator is a straightforward modification of the self-energy
integrals. The 10~K panel (Fig.~8) currently shows artificially sharp
features; lifetime broadening would give the correct low-$T$ linewidth.

\item \textbf{DFT-derived VOC vibrational spectra.}
The molecular database (\code{config/molecular\_database.py}) already
contains 10 DFT-derived volatile organic compounds (benzene, thiophene,
naphthalene, dimethyl sulfide, acetone, etc.) with realistic multi-mode
vibrational spectra. These are deferred until the methodology paper is
submitted because they add a confound between molecule-geometry and
IETS-signature that complicates the message. The rank-1 theorem handles
multiple modes naturally (each shares the same $|u\rangle$), so the
code extension is minimal---only the per-mode $D_0^2$ calibration needs
DFT input.

\item \textbf{Enantioselectivity (Pandey-line extension).} Exploit the
direction-dependent coupling of chiral vibrations to the transport axis to
derive an enantiomer-discrimination figure of merit for the RTD sensor.
Pandey's thesis showed that chirality-dependent projection of molecular
vibrations onto the transport direction produces distinct $d^2I/dV^2$
signatures. This is a compelling long-term goal because no existing
electronic-nose technology can distinguish enantiomers.

\item \textbf{Spectrum-to-percept mapping.} Integrate Pandey's
graph-neural-net mapping of vibrational spectra to olfactory percepts;
demonstrate that device-level $\Delta D(V)$ signals predict perceptual
categories on a held-out test set of known odorants. This connects the
physics simulation to olfactory neuroscience and would be a high-impact
interdisciplinary result.

\item \textbf{Ensemble averaging of molecular position.}
The SISPAD results report single-molecule, optimal-placement response.
The ensemble-averaged response (molecule placed at random transverse
positions over the pixel, with random orientations) is
$\sim\!\tfrac12$ of the optimal value by rough estimate. Computing
this distribution populates the full paper's uncertainty analysis.

\item \textbf{Noise-floor quantification and SNR.}
A realistic $1/f$ contact-fluctuation model and a lock-in detector of
stated time constant would allow $\Delta D(V)$ to be reported as a
signal-to-noise ratio. This is essential for assessing practical
detectability: whether the $\sim$2~nA $\Delta I$ and
$\sim$1~$\mu$A/V$^2$ $\Delta D$ signals are measurable in a real
$10\,\mu$m$\times 10\,\mu$m mesa at room temperature.

\item \textbf{Array-level discrimination.} A pixel-array
pattern-recognition stage (linear SVM or random forest) on a panel of
analytes, to quantify array-level classification accuracy at 300~K.
This extends the single-device question ``can it discriminate?'' to the
system-level question ``how well does an e-nose array classify?''

\item \textbf{Experimental collaboration.} The ZnO/Mg$_{0.3}$Zn$_{0.7}$O
stack is fabricatable today in at least two collaborator labs. A
co-designed experimental demonstration---fabrication, mesa patterning,
lock-in $d^2I/dV^2$ at 300~K---is the natural next proposal cycle.
Krishnamoorthy~2002 demonstrated ZnO/MgZnO RTD NDR; the extension to
IETS requires only depositing analyte molecules on the emitter surface
before capping.

\item \textbf{Alternative filter architectures.} Revisit quantum dots as
the filter element for cases where the RTD geometry cannot be shrunk to
molecular scale. QDs offer stronger confinement and potentially sharper
resonances but are harder to fabricate reproducibly.
\end{enumerate}

\subsection{Mapping sections of this document to target outputs}
\begin{center}
\small
\begin{tabular}{lcccc}
\toprule
Section & SISPAD & TED paper & Thesis chapter & Onboarding doc\\
\midrule
Intro/Motivation (\S1)     & trim to 1 col    & Chap.~1 Intro         & Chap.~1    & \checkmark\\
Lineage (\S2)              & cite only        & Chap.~1 / related work & Chap.~2   & \checkmark\\
Framework (\S3)            & one equation      & Chap.~2 full derivation & Chap.~3  & \checkmark\\
Rank-1 theorem (\S4)       & boxed result \& one para & Chap.~2 + appendix & Chap.~3 & --- \\
Stack selection (\S5)      & 2 sentences       & Chap.~3 full           & Chap.~4   & \checkmark\\
Methodology (\S6)          & grid table only   & Chap.~3 full           & Chap.~4   & \checkmark\\
Validation (\S7)           & one figure        & Chap.~4 full           & Chap.~5   & \checkmark\\
Results (\S8)              & Figs 3--5         & Chap.~5 full           & Chap.~6   & --- \\
Limitations (\S9)          & final paragraph   & Chap.~6 Discussion     & Chap.~7   & \checkmark\\
Roadmap (\S10)             & last paragraph    & Chap.~7 Conclusion     & Chap.~8   & \checkmark\\
\bottomrule
\end{tabular}
\end{center}

% =====================================================================
\appendix
\section{Patil reference parameters (for reproducibility)}
% =====================================================================
\small
\begin{tabular}{ll}
\toprule
Parameter & Value \\
\midrule
Device         & 1D symmetric double-barrier RTD, Patil 2018 \\
Barriers       & Al$_{0.3}$Ga$_{0.7}$As, $\Delta E_c = 0.23$~eV, 2.0~nm each \\
Well           & GaAs, 5.0~nm \\
$m^*$          & $0.067\,m_0$ \\
$a$ (lattice)  & 0.1~nm \\
$E_F$          & 20~meV \\
Bias window    & 0--0.5~V (symmetric: $\mu_{L,R} = E_F \pm V/2$) \\
Modes used     & 2 Einstein modes, $\hbar\omega = 36$, 120~meV \\
Coupling       & $D_0^2 = 0.1$~eV$^2$ for both modes \\
$T$            & 300~K \\
$dE$           & 0.5~meV; $\eta = 1$~meV \\
Our match      & NDR peak 0.55\% at $V = 0.336$~V \\
\bottomrule
\end{tabular}

\section{Bug-fix log (relevant to current code state)}
\begin{enumerate}[leftmargin=1.3em,itemsep=0pt]
\item \textbf{Fix 1 (2026-03):} FBA final-pass $G$ recomputation. Was
producing identical currents for all molecules because $G$ was never
recomputed with the updated phonon $\Sigma$ before the current integral.
\item \textbf{Fix 2:} Coupling normalisation. Patil's Dnu\,$=\,0.1$ is
$D^2$, not $D$. Our \code{self\_energy.py} uses $D^2$ directly; the
molecular database was rescaled by $\sqrt{31.6}$ to match.
\item \textbf{Fix 3:} Contact in-scattering on the inelastic mode.
Missing $f_L\Gamma_L + f_R\Gamma_R$ on the phonon-dressed modes was the
cause of the ``all molecules look identical'' artefact.
\item \textbf{Fix 4:} Retarded phonon self-energy sign. Enforced the
causal form $\Sigma^R_{\mathrm{ph}} = -\tfrac{i}{2}(\Sigma^> - \Sigma^<)$.
\item \textbf{Fix 5:} Landauer shortcut for coherent modes (pre-rank-1).
\item \textbf{Fix 6:} Longitudinal vibron footprint (superseded by the
Gaussian $\chi(z)$ of the rank-1 deriv
ation).
\item \textbf{Fix 7 (2026-04-11):} Baseline $\chi$ restricted to active region only.
The bulk phonon form factor was previously $\chi \equiv 1$ everywhere including contacts,
causing unphysical contact-phonon scattering. Fixed to $\chi_{\mathrm{bulk}} = 1$
only in the barrier+well region (sites where $U_B > 0$), zero in contacts.
\item \textbf{Fix 8 (2026-04-13):} Per-mode $\chi$ vectors. The
\code{fba\_phonon\_sigma} function was applying the same $\chi$ to all phonon
modes. For mixed bulk+molecular simulations, each mode now receives its own
$\chi$: bulk gets $\chi_{\mathrm{bulk}}$ (active region), molecular modes
get the Gaussian $\chi(z)$.
\item \textbf{Fix 9 (2026-04-13):} Bulk coupling calibration. Original
$D^2_{\mathrm{bulk}} = 0.1$~eV$^2$/site over 35 active sites caused 7--9\%
conservation residual and SCBA non-convergence. Root cause: total coupling
$D^2 \times N_{\mathrm{sites}} = 3.5$~eV$^2$ overwhelmed the self-energy.
Reduced to $D^2_{\mathrm{bulk}} = 0.005$~eV$^2$/site so that
$0.005 \times 35 \approx 0.1 \times 3$ (matching the molecular vibron's
effective coupling weight). Conservation improved to $<1$\%.
\item \textbf{Fix 10 (2026-04-14):} SCBA convergence criterion changed from
absolute ($\sum|\Delta\Sigma|$) to relative ($\|\Delta\Sigma\|/\|\Sigma\|$).
The old absolute sum over $\sim 6.8 \times 10^6$ matrix elements made
$\mathrm{tol} = 10^{-6}$ unreachable. With the relative criterion,
convergence at $\mathrm{tol} = 10^{-5}$ is achieved for $\sim$60\% of bias
points within 30 iterations; the remainder (mainly near NDR) still oscillate.
This was a stepping stone; Anderson mixing (Fix~11) then solved the problem
completely.
\item \textbf{Fix 11 (2026-04-14):} Anderson mixing (DIIS) for SCBA convergence.
Stores last $M = 8$ input/residual pairs; solves constrained least-squares via
Lagrange multiplier. With damping $\beta = 0.3$ and Tikhonov regularisation
$10^{-12}$, all bias points now converge in 10--60 iterations, including the
NDR region. Also: $D^2_{\mathrm{bulk}}$ further reduced from 0.005 to 0.001
to eliminate FBA-era near-singular spikes. SCBA with Anderson mixing became
the production solver, superseding FBA. Also added \code{\_robust\_solve()}
wrapper that catches \code{LinAlgError} and adds $10^{-8}j$ regularisation.
Also added ``FINAL PASS'' to the rank-1 solver (\code{scba\_rank1\_keldysh.py})
--- recomputes $G$ with converged $\Sigma_{\mathrm{ph}}$ before the current
integral (the rank-1 analogue of Fix~1).
\item \textbf{Perf 1:} Vectorised $G^R$ solve (\S\ref{sec:methodology}).
\item \textbf{Perf 2:} \code{np.einsum}$\to$\code{np.matmul} on
$(G^R,\Sigma,G^A)$, $\sim$100$\times$ speedup on NumPy 2.4 / Python 3.14.
\end{enumerate}

\section{Glossary}
\begin{description}[leftmargin=1.5em,style=nextline,itemsep=0pt]
\item[IETS] Inelastic Electron Tunneling Spectroscopy.
\item[RTD] Resonant Tunneling Diode.
\item[NEGF] Non-Equilibrium Green's Function.
\item[SCBA] Self-Consistent Born Approximation.
\item[FBA] First Born Approximation, here = SCBA with \code{max\_iter = 1}
and final-pass $G$ recomputation.
\item[NDR] Negative Differential Resistance.
\item[CBO] Conduction-Band Offset.
\item[BEOL] Back-End-Of-Line (CMOS integration).
\item[Fix-$A$] Datta's analytic transverse integration; the transverse area
$A$ enters only as a current prefactor.
\item[Rank-1 projected SCBA] The core methodology of this project; the
full coupled 3D SCBA is reduced to a single 1D Dyson equation because the
vibron coupling matrix in the transverse-mode basis is rank-1.
\end{description}

% =====================================================================
\section*{References}
\addcontentsline{toc}{section}{References}
% =====================================================================
%% Every reference below has been verified against its DOI, the Semantic
%% Scholar API, or the local PDF.  Entries marked NOTE need one remaining
%% detail confirmed before submission.

\begin{thebibliography}{99}

\bibitem{jaklevic1966} R.~C.~Jaklevic and J.~Lambe, ``Molecular vibration
spectra by electron tunneling,'' \textit{Phys.\ Rev.\ Lett.}, vol.~17,
pp.~1139--1140, 1966. DOI: 10.1103/PhysRevLett.17.1139.

\bibitem{patil2018} A.~Patil, D.~Saha, and S.~Ganguly, ``A quantum
biomimetic electronic nose sensor,'' \textit{Sci.\ Rep.}, vol.~8, 128,
2018. DOI: 10.1038/s41598-017-18346-2.

\bibitem{datta} S.~Datta, \textit{Quantum Transport: Atom to Transistor}.
Cambridge University Press, 2005.

\bibitem{meir1992} Y.~Meir and N.~S.~Wingreen, ``Landauer formula for the
current through an interacting electron region,'' \textit{Phys.\ Rev.\
Lett.}, vol.~68, pp.~2512--2515, 1992.

\bibitem{galperin2007} M.~Galperin, M.~A.~Ratner, and A.~Nitzan,
``Molecular transport junctions: vibrational effects,''
\textit{J.\ Phys.: Condens.\ Matter}, vol.~19, 103201, 2007.

\bibitem{troisi2006} A.~Troisi and M.~A.~Ratner, ``Molecular transport
junctions: propensity rules for inelastic electron tunneling spectra,''
\textit{Nano Lett.}, vol.~6, no.~8, pp.~1784--1788, 2006.

\bibitem{tampo2019} \textcolor{red}{\textbf{[REMOVED --- fabricated author list.}
The DOI 10.1109/NANO46743.2019.8874570 exists but with different authors than
``Tampo, Fons, Yamada, Matsubara.'' This reference has been replaced by
Krishnamoorthy~2002 throughout.\textbf{]}}

\bibitem{mgzno_cbo} H.~Yin, J.~Chen, Y.~Wang, J.~Wang, and H.~Guo,
``Composition dependent band offsets of ZnO and its ternary alloys,''
\textit{Sci.\ Rep.}, vol.~7, 41567, 2017. DOI: 10.1038/srep41567.

\bibitem{in2o3kappa} Y.~Kuang, X.~Chen, T.~Ma, Q.~Du, Y.~Zhang, J.~Hao,
F.~Ren, B.~Liu, S.~Zhu, S.~Gu, R.~Zhang, Y.~Zheng, and J.~Ye, ``Band
alignment and enhanced interfacial conductivity manipulated by polarization
in a surfactant-mediated grown $\kappa$-Ga$_2$O$_3$/In$_2$O$_3$
heterostructure,'' \textit{ACS Appl.\ Electron.\ Mater.}, vol.~3,
pp.~795--803, 2021. DOI: 10.1021/acsaelm.0c00947.

\bibitem{lake1997} R.~Lake, G.~Klimeck, R.~C.~Bowen, and D.~Jovanovic,
``Single and multiband modeling of quantum electron transport through
layered semiconductor devices,'' \textit{J.\ Appl.\ Phys.}, vol.~81,
pp.~7845--7869, 1997.

\bibitem{pasupathy2005} A.~N.~Pasupathy \textit{et al.},
``Vibration-assisted electron tunneling in C$_{140}$ transistors,''
\textit{Nano Lett.}, vol.~5, no.~2, pp.~203--207, 2005.
%% Local PDF: papers/C140.pdf. Relevant experimental IETS context.

\bibitem{krishnamoorthy2002} S.~Krishnamoorthy, A.~A.~Iliadis,
A.~Inumpudi, S.~Choopun, R.~D.~Vispute, and T.~Venkatesan,
``Observation of resonant tunneling action in
ZnO/Zn$_{0.8}$Mg$_{0.2}$O devices,'' \textit{Solid-State Electron.},
vol.~46, pp.~1633--1637, 2002. DOI: 10.1016/S0038-1101(02)00117-X.
%% Uses x=0.2; our simulation uses x=0.3 (higher barrier, aids confinement).

\bibitem{sollner1983} T.~C.~L.~G.~Sollner, W.~D.~Goodhue,
P.~E.~Tannenwald, C.~D.~Parker, and D.~D.~Peck, ``Resonant tunneling
through quantum wells at frequencies up to 2.5~THz,''
\textit{Appl.\ Phys.\ Lett.}, vol.~43, pp.~588--590, 1983.
DOI: 10.1063/1.94434.

\bibitem{anderson1965} D.~G.~Anderson, ``Iterative procedures for
nonlinear integral equations,'' \textit{J.\ ACM}, vol.~12, no.~4,
pp.~547--560, 1965.

\bibitem{bayram2010} C.~Bayram, Z.~Vashaei, and M.~Razeghi,
``Room-temperature negative differential resistance characteristics of
GaN/AlN double-barrier resonant tunneling diodes,''
\textit{Appl.\ Phys.\ Lett.}, vol.~96, 042103, 2010.

\end{thebibliography}

\end{document}
